% ECONOMICS_PROBLEM: {"id": "G11-397", "title": "Home bias in cross-border portfolios", "jel_code": "G11", "source_id": "OP-0338"}
% !TeX program = xelatex
\documentclass[11pt,a4paper]{article}
\usepackage[margin=23mm]{geometry}
\usepackage{fontspec}
\setmainfont{Latin Modern Roman}
\usepackage{amsmath,amssymb,mathtools}
\usepackage{xcolor,enumitem,booktabs,longtable,array,xurl,hyperref}
\definecolor{muted}{HTML}{53636C}
\hypersetup{colorlinks=true,urlcolor=blue,linkcolor=blue}
\setlength{\parindent}{0pt}
\setlength{\parskip}{5pt}
\newcommand{\R}{\mathbb R}
\newcommand{\E}{\mathbb E}
\newcommand{\N}{\mathbb N}
\newcommand{\ind}{\mathbf 1}
\newcommand{\Var}{\operatorname{Var}}
\newcommand{\Cov}{\operatorname{Cov}}
\newcommand{\supp}{\operatorname{supp}}
\DeclareMathOperator*{\argmin}{arg\,min}
\DeclareMathOperator*{\argmax}{arg\,max}
\newcommand{\entrymeta}[3]{{\sffamily\small\color{muted}\textbf{\texttt{#1}}\quad|\quad #2\par #3\par}}
\newcommand{\Problemref}[1]{\texttt{#1}}
\newcommand{\JELref}[2]{\texttt{#2} (JEL \texttt{#1})}
\begin{document}
\section*{Home bias in cross-border portfolios}
\textbf{Problem ID:} \texttt{G11-397}\quad \textbf{JEL:} \texttt{G11}\par
% BEGIN ECONOMICS STATEMENT
\label{problem:OP-0338}
{\sffamily\small\color{muted}Original subject group: International finance and exchange rates\par}
\label{definitions:problem:30007620}
\textbf{Audit classification:} Published research agenda. The 2025 paper explicitly proposes narrower further work in footnote 12 (printed p.20), on bilateral political-risk measures, and footnote 15 (printed p.23), on information linkages. Section II.E (pp.13--14) instead supplies existing alternative microfoundations. These limited directions support the agenda classification; the exact joint identification and causal decomposition of all home-bias mechanisms is editorial, with no unique decomposition or current-unsolved status certified.
\subsection*{Definitions and precise research target}
Fix investors from country \(c\), asset class \(j\) and date \(t\). Let \(w_{cjt}^{\mathrm{foreign}}\) be the observed share allocated to foreign claims and \(m_{cjt}^{\mathrm{foreign}}\) the foreign share of a prespecified investable world-market portfolio. For \(m_{cjt}^{\mathrm{foreign}}>0\), define
\[H_{cjt}=1-w_{cjt}^{\mathrm{foreign}}/m_{cjt}^{\mathrm{foreign}}.\]
This index compares holdings with a benchmark, not automatically with an investor's utility-maximizing allocation. A structural demand model includes domestic-income and exchange-rate hedging needs, taxes and transaction costs, legal restrictions, information-acquisition costs and familiarity-related belief distortions. The task is to identify the causal contribution of each mechanism to \(H_{cjt}\), using quantities, returns, investor characteristics and plausibly exogenous changes in those mechanisms. Define each contribution as the difference under a specified mechanism-removal intervention with other equilibrium responses retained. Address interactions and observational equivalence: different microfoundations may generate the same demand equation. Restate securities by ultimate issuer nationality when tax-haven domicile distorts geographic exposure, and distinguish domestic multinational exposure from directly foreign holdings. Report bounds where beliefs or hedging motives are unobserved. The classic equity phenomenon does not establish that every cross-border asset displays the same bias or that all deviations from world weights are irrational.

\textbf{Additional formal definitions and scope (editorial).} Take \(w^{\rm foreign},m^{\rm foreign}\in[0,1]\), \(m^{\rm foreign}>0\), and define the investable benchmark after legal eligibility and free-float adjustments. The index can be negative, indicating foreign overweight, and lies in \([1-1/m^{\rm foreign},1]\). Country of investor residence, issuer nationality and ultimate operating exposure are distinct geographic concepts; state which maps claims to domestic/foreign. For a mechanism-removal policy \(a_k\), define the contrast \(H(0)-H(a_k)\) under re-solved portfolio demand and asset prices. The contrasts need not add to observed bias because hedging, information and taxes interact. A unique additive allocation requires a declared counterfactual order or averaging convention. The world capitalization portfolio is a descriptive benchmark, not a welfare optimum unless additional asset-pricing and diversification assumptions apply.

For the identification part, declare an admissible class \(\Theta\) of structural models, including preferences, technologies, shock laws, measurement/sampling rules and any equilibrium selector. If \(O\) denotes the complete observable history and \(T(\theta)\) the target defined above, the identified set is \(\mathcal I_T=\{T(\theta):\theta\in\Theta,\ \mathcal L_\theta(O)=\mathcal L_{\rm obs}(O)\}\); point identification means this set is a singleton. A \(do(a)\) comparison replaces stated policy/technology equations while fixing the declared exogenous law and re-solving the remaining equations. These definitions do not assert that an identification theorem holds without further restrictions.
\subsection*{Variants and assumption changes}
\begin{enumerate}
\item \textbf{Equity capitalization benchmark (source-motivated).} Measure foreign equity underweight using restated ultimate issuer exposures and explicitly defined investability.
\item \textbf{Mechanism identification across assets (editorial).} Include taxes, hedging and beliefs, separately varying their primitives; test whether the resulting home-bias contrasts are point identified or interactive sets.
\end{enumerate}
\subsection*{References and source audit}
\begin{enumerate}
\item Author-hosted AER scan verifies French and Poterba, May 1991, 81(2):222--226; no DOI assigned. Locator: Opening, p.222; Section III, pp.224--225. Support: Historical equity home-bias background. Full PDF accessible. URL: \url{https://economics.mit.edu/sites/default/files/publications/2006858.pdf}.
\item Author-hosted published QJE PDF verifies Pellegrino, Spolaore and Wacziarg, 2025 and DOI. Locator: Footnote12, printed p.20, and footnote15, printed p.23. Support: Limited further work on bilateral political risk and information linkages. Section II.E, pp.13--14, and the abstract are background/results; the exact joint mechanism-identification question is editorial. Full PDF accessible. URL: \url{https://www.anderson.ucla.edu/faculty_pages/romain.wacziarg/downloads/2025_GCA.pdf}.
\end{enumerate}
\begin{enumerate}\raggedright
\item\hypertarget{definitions:source:30007620:1}{} Kenneth R. French; James M. Poterba. 1991. \emph{Investor Diversification and International Equity Markets}. Opening, p.222; Section III, pp.224--225.
\url{https://economics.mit.edu/sites/default/files/publications/2006858.pdf}.
\item\hypertarget{definitions:source:30007620:2}{} Bruno Pellegrino; Enrico Spolaore; Romain Wacziarg. 2025. \emph{Barriers to Global Capital Allocation}. Footnote12, printed p.20; footnote15, printed p.23 (limited further work); Section II.E, pp.13--14 and abstract (background/results).
\url{https://www.anderson.ucla.edu/faculty_pages/romain.wacziarg/downloads/2025_GCA.pdf}.
DOI: \url{https://doi.org/10.1093/qje/qjaf031}.
\end{enumerate}

\subsection*{Common conventions: Source mathematical conventions}

These conventions apply to each entry unless explicitly replaced.

\textbf{Models and identification.} A primitive model \(m\) specifies all probability laws, feasible choices, information, equilibrium and measurement rules used by its target. The admissible class \(\mathcal M\) is fixed independently of outcomes. If \(P_O(m)\) is its observable probability law and \(\tau(m)\) the target, its identified set at observable law \(P\) is
\[
 \mathcal I_\tau(P)=\{\tau(m):m\in\mathcal M,\ P_O(m)=P\}.
\]
Point identification means that this set is a singleton. An empty set signals incompatibility of data and assumptions. Coordinate bounds need not describe the joint set. Bounds are sharp only if every retained target is attainable or arbitrarily approachable by compatible models. Finite bounds require explicit restrictions on unobserved quantities; unrestricted confounding, hidden wealth, extrapolation or arbitrary equilibrium selection can yield uninformative sets.

\textbf{Causal interventions.} A potential outcome \(Y_i(p)\) is the outcome under a complete policy regime \(p\), including timing, financing, information and market spillovers. The target population and units are fixed before treatment. With interference, use the complete assignment vector or a justified exposure mapping; individual no-interference notation alone cannot identify an economy-wide intervention. A difference of mean potential outcomes does not require the cross-world joint distribution. Conditioning on post-treatment migration, survival, enrollment or employment changes the estimand and needs separate justification.

\textbf{Statistical statements.} An estimator, confidence set, forecast or test specifies a sampling law, model class and loss/coverage criterion. Uniform \((1-\alpha)\)-coverage means \(\inf_{m\in\mathcal M}\Pr_m\{\tau(m)\in C_n\}\geq1-\alpha\), or an explicitly asymptotic version. The whole parameter space trivially has coverage, so informativeness requires a stated width, risk or power objective. A forecasting improvement is relative to a named benchmark, information set and evaluation population.

\textbf{Equilibrium and optimization.} An equilibrium comprises feasible plans, optimality, beliefs where needed, budget constraints and all market/resource clearing. The primitives or a stated benchmark define each correspondence \(\mathcal E(p;m)\). Nonemptiness is assumed or proved, never inferred just from compact policy sets. A selected-equilibrium target uses a declared selection rule; a robust target quantifies over all permitted equilibria. Use suprema or infima unless attainment is established. An \(\varepsilon\)-optimal policy is feasible and within \(\varepsilon\) of the finite optimum value. Report an empty feasible set separately.

\textbf{Welfare and accounting.} Normative utility, interpersonal normalization, social weights, numeraire, discounting and the use of public receipts are primitives. Behavioral utility need not coincide with the welfare criterion. Household welfare includes ownership income and rents once; adding profits again to compensating variation generally double-counts them. A monetary equivalent is defined by an explicit transfer and price convention, with monotonicity and range conditions. Average effects, mechanism contributions, transfers and welfare losses are different targets. Infinite sums require convergence and dynamic feasible plans require terminal or no-Ponzi conditions.

\textbf{Source versions.} A working-paper date, revision date and journal publication year are distinguished. A DOI for a journal version is not automatically the DOI of an earlier manuscript. Locators refer to the stated version and printed pages where specified. A metadata-only check does not verify a claim about a full-text conclusion. Every entry's source subsection narrows the provenance claim accordingly.
% END ECONOMICS STATEMENT
\end{document}
